\chapter{Conclusions and Future Work}
\label{conclusions}

This thesis presented CitySense, an Android-first civic reporting platform designed to simplify pothole reporting through mobile capture, offline report storage, backend automation, and geospatial duplicate handling. The final system combines a Flutter mobile application, a local phone-side report queue, a FastAPI backend, a PostgreSQL/PostGIS data layer, and a web dashboard for operational review. Instead of treating each report as an isolated ticket, the system aggregates nearby reports into a single issue with an increasing upvote count. This makes the resulting data more useful for prioritization and more readable on a map.

The work demonstrates that a practical civic reporting pipeline can be built by integrating established technologies rather than inventing a new algorithm from scratch. Flutter supports a simple citizen workflow, SQLite keeps captured reports safe until synchronization, FastAPI and SQLAlchemy provide a clear service boundary, PostGIS supplies the spatial reasoning needed for duplicate merging, and a YOLO-based detector enables automation of the issue classification step. Together, these parts form a coherent software prototype with a clear civic reporting purpose.

At the same time, the project remains intentionally scoped. The system focuses on potholes rather than on a wide taxonomy of municipal issues. A YOLOv8n detector was fine-tuned on a pothole dataset, integrated into the backend as \texttt{pothole\_yolov8n.pt}, evaluated on validation and test splits, and configured with an application confidence threshold of 0.65 after calibration and field testing. This boundary is important because it keeps the written document aligned with the implemented evidence: CitySense is a working pothole-first prototype, not yet a general urban-damage detection platform.

\section{Future Work}

Several future directions follow naturally from the current version. A next version could collect and annotate additional classes such as trash piles, damaged signs, and damaged pole-mounted trash bins, then retrain the detector as a multi-class YOLO model and update the backend label mapping accordingly. It could also add user accounts, municipal roles, and notification workflows. Another useful extension would be to improve offline synchronization with background scheduling and clearer conflict-resolution tools, connect the dashboard to maintenance scheduling or work-order systems, and improve map analytics through clustering, heatmaps, and historical trends.

In summary, CitySense should be understood as a strong prototype and a foundation for subsequent development. Its main accomplishment is not only that it detects or stores potholes, but that it demonstrates a complete reporting loop in which image capture, location, geospatial reasoning, and administrative review are treated as parts of one system.
